\documentclass[11pt,a4paper]{article}
\input{casper_preambule}
\rhead{\small Casper --- Corrigé}
\renewcommand{\Q}{\medskip\par\stepcounter{qcnt}\noindent\textbf{\color{cnc}Question~\theqcnt.}~}
\begin{document}
\begin{center}
{\large\bfseries Préparation au Concours National Commun --- Informatique MP}\\[8pt]
\rule{0.9\linewidth}{0.5pt}\\[6pt]
{\Large\bfseries Corrigé détaillé}\\[4pt]
{\bfseries Filière MP --- Durée : 4 heures}\\[6pt]
{\itshape Finalité dans les chaînes de blocs : le protocole Casper}\\[2pt]
{\small d'après V.~Buterin et V.~Griffith, \emph{Casper the Friendly Finality Gadget}, \\ Ethereum Foundation, arXiv:1710.09437v4 (2019)}\\[6pt]
\rule{0.9\linewidth}{0.5pt}\\[8pt]
{\small Auteur~: \href{https://orcid.org/0000-0002-6108-7176}{\textbf{Pr Agrégé El Hadiq Zouhair}}}\\[3pt]
{\small \href{https://cpgeacademy.org}{\textbf{cpgeacademy.org}} \quad---\quad Tous droits réservés \textcopyright\ cpgeacademy.org}
\end{center}
\medskip
\begin{center}
\fcolorbox{gray}{gray!5}{\begin{minipage}{0.93\linewidth}\small
\textbf{Avis important.} Ce document est une \textbf{initiative personnelle} du professeur,
à but strictement pédagogique. Il ne s'agit \textbf{pas} d'un sujet officiel du Concours
National Commun~; il s'en inspire uniquement par la forme et l'esprit.
\smallskip

\textbf{Pour aller plus loin.} Les élèves peuvent traiter une \textbf{nouvelle version} de
ce problème dotée d'un \textbf{autocorrecteur des réponses}, et suivre un \textbf{cours
complet}, sur \href{https://cpgeacademy.org}{\textbf{cpgeacademy.org}}.
\smallskip

\textbf{Note.} L'intégralité du code Python de ce corrigé a été exécutée et validée~; les
fonctions de la partie IV et la programmation dynamique de la partie V ont en outre été
confrontées à une recherche exhaustive sur plusieurs milliers de jeux d'essai aléatoires.
\end{minipage}}
\end{center}
\medskip

%=======================================================================
\section*{Partie I --- Points de contrôle, dépôts et fuite d'inactivité}
%=======================================================================

\Q Les deux fonctions se réduisent à une opération arithmétique.
\begin{lstlisting}
def est_checkpoint(n):
    return n % 100 == 0

def hauteur_bloc(n):
    return n // 100
\end{lstlisting}
Chacune effectue une unique division euclidienne~: la complexité est $O(1)$ (en supposant,
comme d'usage, le coût des opérations arithmétiques sur les entiers manipulés constant).

\Q On remonte les liens de parenté jusqu'à la racine.
\begin{lstlisting}
def hauteur(parent, c):
    h = 0
    x = c
    while parent[x] != -1:
        x = parent[x]
        h = h + 1
    return h
\end{lstlisting}
\textbf{Terminaison.} Posons $V = x$ (la valeur courante de la variable \texttt{x}). C'est un
entier naturel. À chaque tour de boucle, $x$ est remplacé par \texttt{parent[x]}~; or
l'hypothèse $0\le\texttt{parent[c]}<c$ garantit $\texttt{parent[x]}<x$. Le variant $V$ est
donc un entier naturel \emph{strictement décroissant}~: la boucle effectue un nombre fini de
tours. \hfill$\square$

\textbf{Complexité.} À chaque tour, la hauteur de $x$ diminue exactement de $1$~; partant de
$h(c)$, la boucle effectue exactement $h(c)$ tours, chacun en $O(1)$. La complexité est donc
$\Theta\big(h(c)\big)$, soit $O(n)$ dans le pire des cas (arbre réduit à un chemin).

\Q L'hypothèse $\texttt{parent[c]}<c$ signifie que les n\oe uds sont numérotés dans un ordre
compatible avec la relation de descendance~: le père est toujours \emph{déjà} traité.
\begin{lstlisting}
def toutes_hauteurs(parent):
    n = len(parent)
    h = [0] * n
    for c in range(n):
        if parent[c] == -1:
            h[c] = 0
        else:
            h[c] = h[parent[c]] + 1
    return h
\end{lstlisting}
\textbf{Invariant.} \emph{Au début de l'itération d'indice $c$, pour tout $c'<c$ on a
$\texttt{h[c']}=h(c')$.}

\emph{Initialisation.} Pour $c=0$ l'ensemble des indices $c'<0$ est vide~: l'invariant est
trivialement vrai.

\emph{Hérédité.} Supposons l'invariant vrai au début de l'itération $c$. Si $c=0$, alors
$c$ est la racine et $h(0)=0$~: l'affectation est correcte. Sinon, $p=\texttt{parent[c]}$
vérifie $0\le p<c$, donc par hypothèse d'induction $\texttt{h[p]}=h(p)$~; comme le chemin de
$c$ à la racine est le chemin de $p$ à la racine augmenté de l'arête $c\to p$, on a
$h(c)=h(p)+1$, ce qu'affecte l'algorithme. L'invariant est vrai au début de l'itération
$c+1$. \hfill$\square$

\emph{Conclusion.} Après la dernière itération, $\texttt{h[c']}=h(c')$ pour tout $c'<n$.

C'est \emph{exactement ici} que l'hypothèse $\texttt{parent[c]}<c$ est indispensable~: sans
elle, \texttt{h[parent[c]]} pourrait ne pas encore avoir été calculée et vaudrait la valeur
d'initialisation $0$, produisant un résultat faux.

\textbf{Complexité.} $\Theta(n)$ en temps et $\Theta(n)$ en espace. À comparer aux
$O(n\cdot h)$ qu'aurait coûté $n$ appels à \texttt{hauteur}.

\Q Avec $\texttt{parent} = [-1,0,1,2,1,4,0,6]$~:
\[
\begin{array}{c|cccccccc}
c & 0&1&2&3&4&5&6&7\\\hline
\texttt{parent[c]} & -1&0&1&2&1&4&0&6\\
h(c) & 0&1&2&3&2&3&1&2
\end{array}
\]
On retrouve bien $h=[0,1,2,3,2,3,1,2]$.

\Q Le point délicat est d'éviter tout flottant.
\begin{lstlisting}
def poids(depots, E):
    s = 0
    for v in E:
        s = s + depots[v]
    return s

def est_supermajorite(depots, E, T):
    return 3 * poids(depots, E) >= 2 * T
\end{lstlisting}
\textbf{Justification.} Comme $3>0$, on a l'équivalence
$w(E)\ge \tfrac{2}{3}T \iff 3\,w(E)\ge 2\,T$. Les deux membres de droite sont des
\emph{entiers}~: le test est exact, alors qu'un test \code{w >= 2*T/3} en flottants pourrait
renvoyer un résultat erroné par arrondi (typiquement lorsque $2T$ n'est pas divisible par
$3$, ou pour des dépôts de l'ordre de $10^{18}$ dépassant la précision des \emph{doubles}).
Complexité~: $O(|E|)$.

\Q $T = 30+30+20+20 = 100$, donc $2T = 200$.
\begin{itemize}[leftmargin=1.4em,itemsep=1pt]
\item $\{a,b\}$~: $w=60$, $3w=180<200$ $\Rightarrow$ \textbf{pas} une supermajorité.
\item $\{a,b,c\}$~: $w=80$, $3w=240\ge 200$ $\Rightarrow$ supermajorité.
\item $\{a,c,d\}$~: $w=70$, $3w=210\ge 200$ $\Rightarrow$ supermajorité.
\end{itemize}

\Q \textbf{(Lemme d'intersection.)} Soient $E_1,E_2$ deux supermajorités.

Le dépôt étant additif sur les ensembles disjoints, la formule d'inclusion-exclusion donne
\[ w(E_1\cup E_2)\;=\;w(E_1)+w(E_2)-w(E_1\cap E_2). \]
Or $E_1\cup E_2$ est inclus dans l'ensemble de tous les validateurs et les dépôts sont
strictement positifs, donc $w(E_1\cup E_2)\le T$. Il vient
\[ w(E_1\cap E_2)\;=\;w(E_1)+w(E_2)-w(E_1\cup E_2)\;\ge\;\frac{2}{3}T+\frac{2}{3}T-T
\;=\;\frac{T}{3}. \]
\hfill$\square$

\emph{Interprétation.} Deux supermajorités quelconques se recoupent toujours sur au moins un
tiers du dépôt total. C'est ce tiers qui sera détruit en cas d'attaque~: toute la sûreté de
Casper repose sur ce calcul élémentaire.

\Q À chaque époque le dépôt est multiplié par $1-p$, donc, par récurrence immédiate,
\[ D_k \;=\; D\,(1-p)^k. \]
La perte subie à l'époque $k$ vaut $p\,D_k$. Comme $0<1-p<1$, la série géométrique converge
et
\[ \sum_{k=0}^{+\infty} p\,D_k \;=\; p\,D\sum_{k=0}^{+\infty}(1-p)^k
\;=\; \frac{p\,D}{1-(1-p)} \;=\; D. \]
\textbf{Interprétation.} La somme des pertes cumulées converge exactement vers le dépôt
initial~: la fuite d'inactivité finit par confisquer la \emph{totalité} du dépôt d'un
validateur définitivement inactif, mais jamais en temps fini (il en reste toujours une
fraction $(1-p)^k>0$). Le mécanisme est donc « doux »~: il ne punit pas brutalement une
absence passagère, mais élimine asymptotiquement les validateurs absents, jusqu'à ce que les
validateurs restants redeviennent une supermajorité.

\Q Comme $0<1-p<1$, on a $(1-p)^k\xrightarrow[k\to+\infty]{}0$, donc il existe $k$ tel que
$D(1-p)^k\le S$~: l'ensemble considéré est une partie non vide de $\mathbb{N}$, il admet un
plus petit élément. Pour $k\in\mathbb{N}$~:
\[ D_k\le S \iff (1-p)^k\le \frac{S}{D}
\iff k\,\ln(1-p)\le \ln\!\Big(\frac{S}{D}\Big)
\iff k\ \ge\ \frac{\ln(S/D)}{\ln(1-p)}, \]
la dernière équivalence provenant de $\ln(1-p)<0$ (le sens de l'inégalité s'inverse). Comme
$0<S<D$, le quotient $x=\ln(S/D)/\ln(1-p)$ est strictement positif. L'ensemble des entiers
$k\ge x$ a pour plus petit élément $\lceil x\rceil$, d'où
$k^{\star}=\big\lceil \ln(S/D)/\ln(1-p)\big\rceil$. \hfill$\square$

\Q
\begin{lstlisting}
def nb_epoques(D, p, S):
    k = 0
    d = D
    while d > S:
        d = d * (1 - p)
        k = k + 1
    return k
\end{lstlisting}
\textbf{Terminaison.} D'après la question 9, l'entier $k^{\star}$ existe. Posons
$V = k^{\star}-k$. Tant que la boucle continue on a $d=D_k>S$, donc $k<k^{\star}$ (par
minimalité de $k^{\star}$ et décroissance de $k\mapsto D_k$), donc $V\ge 1>0$~; et $V$
décroît strictement de $1$ à chaque tour. C'est un variant entier positif~: la boucle
termine, et elle s'arrête précisément lorsque $k=k^{\star}$. \hfill$\square$

\textbf{Complexité.} Le nombre de tours est $k^{\star}=\big\lceil
\ln(D/S)/\ln\!\big(\tfrac{1}{1-p}\big)\big\rceil$. À $p$ fixé, $1/\ln\frac{1}{1-p}$ est une
constante, donc $k^\star = O\!\big(\log(D/S)\big)$, chaque tour coûtant $O(1)$.

\Q Pour $D=1000$, $p=0{,}05$, $S=500$~:
$x = \dfrac{\ln(0{,}5)}{\ln(0{,}95)} = \dfrac{-0{,}693147}{-0{,}051293} \approx 13{,}513$,
donc $k^{\star}=\lceil 13{,}513\rceil = \mathbf{14}$.

\emph{Vérification.} $D_{13}=1000\times0{,}95^{13}\approx 513{,}3>500$ et
$D_{14}=1000\times0{,}95^{14}\approx 487{,}7\le 500$. Il faut donc $14$ époques pour qu'un
validateur inactif perde la moitié de son dépôt.

%=======================================================================
\section*{Partie II --- Arbre de points de contrôle et sérialisation des votes}
%=======================================================================

\Q
\begin{lstlisting}
def est_ancetre(parent, a, b):
    x = b
    while x != -1:
        if x == a:
            return True
        x = parent[x]
    return False
\end{lstlisting}
\textbf{Terminaison.} Le variant est $V=x+1\in\mathbb{N}$~: en effet $x\ge -1$ et, à chaque
tour, $x$ est remplacé par $\texttt{parent[x]}<x$. $V$ est un entier naturel strictement
décroissant. \hfill$\square$

\textbf{Complexité.} La boucle parcourt le chemin de $b$ à la racine, soit au plus
$h(b)+1$ n\oe uds~: la complexité est $O\big(h(b)\big)$, donc $O(n)$ au pire.

\Q
\begin{lstlisting}
def conflit(parent, a, b):
    return not est_ancetre(parent, a, b) and not est_ancetre(parent, b, a)
\end{lstlisting}
Sur l'arbre $A_0$~:
\begin{itemize}[leftmargin=1.4em,itemsep=1pt]
\item $(2,4)$~: le chemin de $4$ à la racine est $4,1,0$ (ne contient pas $2$) et celui de
$2$ est $2,1,0$ (ne contient pas $4$)~: \textbf{en conflit}.
\item $(3,5)$~: chemins $3,2,1,0$ et $5,4,1,0$~: \textbf{en conflit}.
\item $(1,3)$~: le chemin de $3$ à la racine est $3,2,1,0$ et contient $1$~: $1$ est ancêtre
de $3$, donc \textbf{pas} de conflit.
\end{itemize}

\Q
\begin{lstlisting}
def ancetre_commun(parent, a, b, h):
    x, y = a, b
    while h[x] > h[y]:      # on remonte le plus profond des deux
        x = parent[x]
    while h[y] > h[x]:
        y = parent[y]
    while x != y:           # meme hauteur : on remonte ensemble
        x = parent[x]
        y = parent[y]
    return x
\end{lstlisting}
\textbf{Complexité.} Les deux premières boucles effectuent respectivement au plus $h(a)$ et
$h(b)$ tours, la troisième au plus $\min\big(h(a),h(b)\big)$ tours. Le total est
$O\big(h(a)+h(b)\big)$, soit $O(n)$ au pire.

\Q \textbf{Lemme admis (énoncé précis).} \emph{Si $x$ et $y$ sont deux ancêtres d'un même
n\oe ud $z$ et si $h(x)=h(y)$, alors $x=y$.} (En effet, les ancêtres de $z$ forment le
chemin de $z$ à la racine, sur lequel la hauteur prend chaque valeur de $0$ à $h(z)$
exactement une fois.)

\textbf{Correction des deux premières boucles.} Elles préservent l'énoncé « $x$ est un
ancêtre de $a$ et $y$ un ancêtre de $b$ », puisqu'on ne fait que remonter des liens de
parenté~; elles préservent aussi « tout ancêtre commun de $a$ et $b$ est un ancêtre de $x$
et de $y$ ». Justifions ce second point pour la première boucle~: si $h(x)>h(y)$ et si $u$
est un ancêtre commun de $a$ et $b$, alors $h(u)\le h(y)<h(x)$, donc $u\ne x$~; comme $u$ et
$x$ sont tous deux ancêtres de $a$, $u$ est un ancêtre de $x$ distinct de $x$, donc un
ancêtre de \texttt{parent[x]}. À la sortie, $h(x)=h(y)$.

\textbf{Invariant de la troisième boucle.}
\begin{quote}\itshape
$x$ est un ancêtre de $a$, $y$ est un ancêtre de $b$, $h(x)=h(y)$, et tout ancêtre commun de
$a$ et $b$ est un ancêtre à la fois de $x$ et de $y$.
\end{quote}
\emph{Initialisation.} Vrai à l'entrée d'après ce qui précède.

\emph{Hérédité.} Si $x\ne y$, soit $u$ un ancêtre commun de $a$ et $b$~; par invariant $u$
est un ancêtre de $x$ et de $y$. Si l'on avait $u=x$, alors $u$ serait un ancêtre de $y$ de
même hauteur que $y$, donc $u=y$ par le lemme, d'où $x=y$~: contradiction. Donc $u\ne x$ et
de même $u\ne y$, si bien que $u$ est un ancêtre de \texttt{parent[x]} et de
\texttt{parent[y]}. Les hauteurs restent égales. L'invariant est préservé.

\emph{Terminaison.} $h(x)$ est un entier naturel qui décroît strictement à chaque tour~;
lorsque $h(x)=0$ on a $x=y=r$ et la boucle s'arrête (au plus tard).

\emph{Conclusion.} À la sortie $x=y$ est un ancêtre commun de $a$ et $b$, et tout ancêtre
commun de $a$ et $b$ est un ancêtre de $x$, donc de hauteur $\le h(x)$. C'est bien le plus
proche ancêtre commun. \hfill$\square$

\Q Notons $\ell=\texttt{ancetre\_commun}(a,b)$.

$(\Leftarrow)$ Par contraposée. Si $a$ et $b$ ne sont pas en conflit, l'un est ancêtre de
l'autre~; disons $a$ ancêtre de $b$. Alors $a$ est un ancêtre commun de $a$ et $b$, donc
$h(a)\le h(\ell)$~; mais $\ell$ est un ancêtre de $a$, donc $h(\ell)\le h(a)$. D'où
$h(\ell)=h(a)$ et, $\ell$ et $a$ étant deux ancêtres de $a$ de même hauteur, $\ell=a$ par le
lemme. Donc $\ell\in\{a,b\}$.

$(\Rightarrow)$ Par contraposée également. Si $\ell=a$, alors $a=\ell$ est un ancêtre de
$b$, donc $a$ et $b$ ne sont pas en conflit~; de même si $\ell=b$. \hfill$\square$

\Q
\begin{lstlisting}
def vote_valide(parent, depots, vote):
    v, s, t = vote
    return (v in depots) and (s != t) and est_ancetre(parent, s, t)
\end{lstlisting}
Complexité~: $O\big(h(t)\big)$ en moyenne (le test \code{v in depots} est en $O(1)$ moyen).

\Q
\begin{lstlisting}
def encode(n):
    octets = []
    while True:
        b = n % 128
        n = n // 128
        if n > 0:
            octets.append(b + 128)
        else:
            octets.append(b)
            return octets
\end{lstlisting}
\[
\texttt{encode(0)} = [0],\quad
\texttt{encode(127)} = [127],\quad
\texttt{encode(128)} = [128,\,1],\quad
\texttt{encode(300)} = [172,\,2].
\]
\emph{Détail pour $300$~:} $300 = 2\times 128 + 44$, donc $a_0=44$, $a_1=2$, et le codage est
$[44+128,\,2]=[172,\,2]$.

\Q Notons $n_0=n$ et $n_{k+1}=\lfloor n_k/128\rfloor$ la valeur de la variable \code{n}
après $k+1$ tours. Une récurrence immédiate donne $n_k=\lfloor n/128^k\rfloor$. La boucle
s'arrête au premier tour où $n_{k}=0$, c'est-à-dire produit exactement
\[ \ell = \min\{k\ge 1 \ :\ \lfloor n/128^{k}\rfloor = 0\}
     = \min\{k\ge 1 \ :\ n<128^{k}\} \]
octets. Si $n=0$ alors $\ell=1$. Si $n\ge 1$, $n<128^k \iff \log_{128} n < k$, et le plus
petit entier $k\ge1$ strictement supérieur à $\log_{128}n$ est
$\lfloor \log_{128} n\rfloor+1$. D'où
\[ \ell = \lfloor \log_{128} n\rfloor + 1 = O(\log n). \]
\textbf{Variant.} La suite $(n_k)$ est une suite d'entiers naturels~; tant que la boucle
continue on a $n_k\ge 1$ donc $n_{k+1}=\lfloor n_k/128\rfloor < n_k$. Le variant $n_k$ est
un entier naturel strictement décroissant~: terminaison. \hfill$\square$

\Q
\begin{lstlisting}
def decode(octets):
    n = 0
    p = 1
    for b in octets:
        n = n + (b % 128) * p
        p = p * 128
    return n
\end{lstlisting}
Complexité $\Theta(\ell)$ où $\ell$ est la longueur de la liste.

\Q On utilise d'abord la propriété suivante, immédiate d'après le code~: pour tout octet
$b_0$ et toute liste $L$,
\begin{equation}\label{eq:dec}
\texttt{decode}([b_0]+L) \;=\; (b_0 \bmod 128) \;+\; 128\times \texttt{decode}(L).
\end{equation}

\textbf{Récurrence sur $\ell$}, le nombre d'octets de \code{encode(n)}.

\emph{Cas $\ell=1$.} Cela équivaut à $n<128$ (question 19). Alors $\texttt{encode}(n)=[n]$ et
$\texttt{decode}([n]) = n \bmod 128 = n$.

\emph{Hérédité.} Soit $\ell\ge 2$ et supposons la propriété vraie pour tous les entiers dont
le codage a $\ell-1$ octets. Soit $n$ tel que $\texttt{encode}(n)$ ait $\ell$ octets, donc
$n\ge 128$. Posons $q=\lfloor n/128\rfloor\ge 1$ et $b=n\bmod 128$. Le premier tour de boucle
empile $b+128$ (car $q>0$) puis poursuit avec $q$~: par définition de l'algorithme,
\[ \texttt{encode}(n) \;=\; [\,b+128\,] \;+\; \texttt{encode}(q), \]
et $\texttt{encode}(q)$ a $\ell-1$ octets. En appliquant \eqref{eq:dec} puis l'hypothèse de
récurrence~:
\[ \texttt{decode}(\texttt{encode}(n)) = \big((b+128)\bmod 128\big) + 128\times
\texttt{decode}(\texttt{encode}(q)) = b + 128\,q = n, \]
la dernière égalité étant la division euclidienne de $n$ par $128$. \hfill$\square$

\Q On a $128^4 = 2^{28} = 268\,435\,456 < 10^9$ et $128^5 = 2^{35} = 34\,359\,738\,368
> 10^9$. D'après la question 19, le codage d'une hauteur $h\le 10^9$ occupe donc au plus
$\lfloor\log_{128}10^9\rfloor+1 = \mathbf{5}$ octets.

\textbf{Commentaire.} Un codage sur $8$ octets (entier $64$ bits) consommerait toujours $8$
octets. En pratique les hauteurs de points de contrôle réelles se comptent en milliers~:
$1000<128^2=16\,384$, donc $2$ octets suffisent, soit un gain de $75\,\%$. Un vote contenant
deux hauteurs, l'économie est de $12$ octets par vote~; avec plusieurs dizaines de milliers
de validateurs votant à chaque époque, cela représente plusieurs centaines de kilo-octets
économisés par époque sur la bande passante du réseau et sur le stockage définitif dans la
chaîne. Le prix à payer est un décodage légèrement plus coûteux et la nécessité de connaître
la longueur du codage \emph{a posteriori} (d'où le bit de continuation).

%=======================================================================
\section*{Partie III --- Justification et finalisation}
%=======================================================================

\Q On agrège les dépôts par couple $(s,t)$, en dédoublonnant au préalable les votes.
\begin{lstlisting}
def liens_supermajorite(votes, depots):
    T = 0
    for v in depots:
        T = T + depots[v]
    poids_lien = {}
    vus = set()
    for (v, s, t) in votes:
        if (v, s, t) not in vus:          # un validateur ne compte qu'une fois
            vus.add((v, s, t))
            if (s, t) in poids_lien:
                poids_lien[(s, t)] = poids_lien[(s, t)] + depots[v]
            else:
                poids_lien[(s, t)] = depots[v]
    L = set()
    for st in poids_lien:
        if 3 * poids_lien[st] >= 2 * T:
            L.add(st)
    return L
\end{lstlisting}
\textbf{Correction.} L'ensemble \code{vus} garantit qu'un même triplet $(\nu,s,t)$ répété
n'est comptabilisé qu'une fois~; ainsi $\texttt{poids\_lien[(s,t)]} = w(E_{s\to t})$ où
$E_{s\to t}$ est l'ensemble des validateurs ayant voté $s\to t$. Le test final est celui de
la question 5.

\textbf{Complexité.} $O(m)$ en moyenne pour la boucle principale ($m$ = nombre de votes,
opérations de dictionnaire et d'ensemble en $O(1)$ moyen), plus $O(|L'|)$ pour la boucle
finale où $|L'|\le m$ est le nombre de couples distincts, plus $O(N)$ pour le calcul de $T$.
Total~: $O(m+N)$ en moyenne.

\Q
\begin{lstlisting}
def enfants_de(parent):
    E = [[] for _ in range(len(parent))]
    for c in range(len(parent)):
        if parent[c] != -1:
            E[parent[c]].append(c)
    return E
\end{lstlisting}
Complexité $\Theta(n)$ en temps et en espace (la somme des longueurs des listes vaut $n-1$).

\emph{Remarque.} On écrira \code{[[] for _ in range(n)]} et surtout \emph{pas}
\code{[[]] * n}, qui créerait $n$ références vers une \textbf{même} liste.

\Q
\begin{lstlisting}
def sous_arbre(E, c):
    r = [c]
    for f in E[c]:
        r = r + sous_arbre(E, f)
    return r
\end{lstlisting}
\textbf{Nombre d'appels.} Notons $A(c)$ le nombre total d'appels engendrés par
\code{sous_arbre(E, c)} (lui compris). On a $A(c) = 1+\sum_{f\in E[c]} A(f)$. Par récurrence
sur la hauteur du sous-arbre, $A(c)$ est le nombre de descendants de $c$. Comme chaque
n\oe ud $\ne r$ figure dans la liste des fils d'exactement un n\oe ud (son père, unique), le
parcours atteint chaque n\oe ud une et une seule fois~: $A(0)=n$. \hfill$\square$

\textbf{Coût des concaténations.} L'opérateur \code{+} sur les listes recopie ses deux
opérandes. Sur un arbre réduit à un \emph{chemin} ($\texttt{parent[c]}=c-1$), le sous-arbre
de $c$ a $n-c$ éléments et l'appel de rang $c$ recopie $n-c$ éléments~: le coût total est
\[ \sum_{c=0}^{n-1}(n-c) \;=\; \frac{n(n+1)}{2} \;=\; \Theta(n^2). \]

\textbf{Variante linéaire.} On accumule dans une même liste, sans jamais recopier.
\begin{lstlisting}
def sous_arbre_acc(E, c, acc=None):
    if acc is None:
        acc = []
    acc.append(c)
    for f in E[c]:
        sous_arbre_acc(E, f, acc)
    return acc
\end{lstlisting}
Chaque n\oe ud donne lieu à exactement un \code{append} en $O(1)$ amorti~: la complexité est
$\Theta(n)$ en temps. L'espace de pile est $O(h)$ où $h$ est la hauteur de l'arbre.

\Q La justification est la \emph{clôture} de $\{r\}$ par les liens de supermajorité~: c'est
l'ensemble des sommets accessibles depuis $r$ dans le graphe orienté $(\,\text{n\oe uds},\,L\,)$.
\begin{lstlisting}
def justifies(parent, liens, racine=0):
    succ = {}                       # liste d'adjacence du graphe des liens
    for (s, t) in liens:
        if s in succ:
            succ[s].append(t)
        else:
            succ[s] = [t]
    J = {racine}
    pile = [racine]
    while pile:
        c = pile.pop()
        if c in succ:
            for t in succ[c]:
                if t not in J:
                    J.add(t)
                    pile.append(t)
    return J
\end{lstlisting}
\textbf{Complexité.} Construction de \code{succ} en $O(|L|)$~; chaque sommet est empilé au
plus une fois (il est ajouté à \code{J} au moment où il est empilé) et chaque arc est examiné
au plus une fois. Total~: $O(n+|L|)$ en moyenne.

\Q Notons $J$ l'ensemble renvoyé et $\mathcal{J}$ l'ensemble des points de contrôle
justifiés au sens de la définition.

$\boxed{J\subseteq\mathcal{J}}$ Montrons par récurrence sur l'ordre d'insertion dans $J$ que
tout élément inséré est justifié. Le premier inséré est $r$, justifié par définition.
Supposons les $k$ premiers insérés justifiés et soit $t$ le $(k+1)$-ème~: il a été inséré
lors du traitement d'un sommet $c$ dépilé, donc déjà inséré, donc justifié par hypothèse de
récurrence, et $(c,t)\in L$, c'est-à-dire $c\to t$. Par définition, $t$ est justifié.

$\boxed{\mathcal{J}\subseteq J}$ Soit $c\in\mathcal{J}$. Par définition, il existe une chaîne
$r=b_0\to b_1\to\dots\to b_k=c$ de liens de supermajorité (chaque $b_{i}$ justifiant
$b_{i+1}$). Montrons par récurrence sur $k$ que $c\in J$. Si $k=0$, $c=r\in J$ par
initialisation. Si $k\ge 1$, par hypothèse de récurrence $b_{k-1}\in J$~; or tout sommet
inséré dans $J$ est empilé, donc dépilé (la boucle ne s'arrête que lorsque la pile est vide,
et elle s'arrête car chaque sommet est empilé au plus une fois). Au moment où $b_{k-1}$ est
dépilé, l'arc $(b_{k-1},c)$ est examiné~: soit $c$ est déjà dans $J$, soit il y est ajouté.
Dans les deux cas $c\in J$. \hfill$\square$

\Q On parcourt directement les liens plutôt que le produit $J\times L$.
\begin{lstlisting}
def finalises(parent, liens, J, h, racine=0):
    F = {racine}
    for (s, t) in liens:
        if s in J and parent[t] == s and h[t] == h[s] + 1:
            F.add(s)
    return F
\end{lstlisting}
\emph{Remarque.} Le test \code{parent[t] == s} implique déjà \code{h[t] == h[s] + 1}~; on
conserve les deux pour coller littéralement à la définition.

\textbf{Complexité.} $O(|L|)$ en moyenne. (La version naïve parcourant $J\times L$ coûterait
$O(|J|\cdot|L|)$.)

\Q Soit $c\ne r$ finalisé. Par définition, $c$ est justifié et il existe un lien $c\to c'$
où $c'$ est un fils direct de $c$. Or la définition de la justification stipule qu'un point
de contrôle $x$ est justifié dès qu'il existe un point justifié $y$ avec $y\to x$. En prenant
$y=c$ (justifié) et $x=c'$, on conclut que $c'$ est justifié. \hfill$\square$

\emph{Conséquence.} Finaliser $c$, c'est justifier deux points de contrôle
\textbf{consécutifs} $c$ et $c'$ avec $h(c')=h(c)+1$~: c'est cette « double justification
contiguë » qui rend impossible tout emboîtement ultérieur, et donc toute révision de $c$.

\Q Le dépôt total vaut $T=100$ et $2T=200$.

\textbf{Liens.} Chacun des couples $(0,1)$, $(1,2)$, $(2,3)$ est voté par $\{a,b,c\}$, de
poids $30+30+20=80$~; or $3\times 80 = 240 \ge 200$. Les trois couples sont donc des liens
de supermajorité, et ce sont les seuls~:
\[ L = \{(0,1),\,(1,2),\,(2,3)\}. \]

\textbf{Justifiés.} $0=r$ est justifié~; $0\to 1$ donne $1$~; $1\to 2$ donne $2$~; $2\to 3$
donne $3$. Aucun autre n\oe ud n'est cible d'un lien. Donc
$J=\{0,1,2,3\}$.

\textbf{Finalisés.} $0=r$ est finalisé par définition. Pour $1$~: il est justifié, le lien
$1\to 2$ existe et $\texttt{parent[2]}=1$ avec $h(2)=h(1)+1=2$ $\Rightarrow$ finalisé. Pour
$2$~: justifié, lien $2\to 3$, $\texttt{parent[3]}=2$, $h(3)=3=h(2)+1$ $\Rightarrow$
finalisé. Pour $3$~: il est justifié mais aucun lien ne part de $3$ $\Rightarrow$ non
finalisé. Donc
\[ F=\{0,1,2\}. \]

\Q
\begin{lstlisting}
def choix_fourche(J, h):
    meilleur = None
    for c in J:
        if meilleur is None or h[c] > h[meilleur]:
            meilleur = c
    return meilleur
\end{lstlisting}
Complexité $\Theta(|J|)$, soit $O(n)$. (L'ensemble $J$ est non vide car il contient toujours
la racine.)

\Q \textbf{(a)} Soit $c$ justifié. Si $c=r$ la chaîne est réduite à $b_0=r$. Sinon, la
définition fournit un justifié $b_{k-1}$ avec $b_{k-1}\to c$~; en itérant le raisonnement sur
$b_{k-1}$, on construit une suite finie $r=b_0\to b_1\to\dots\to b_k=c$ de justifiés (le
procédé termine car, un lien $s\to t$ n'existant que si $s$ est un ancêtre \emph{strict} de
$t$, la suite des hauteurs $h(b_i)$ est strictement décroissante quand on remonte, donc
atteint $0$, c'est-à-dire la racine, en un nombre fini d'étapes).

Par définition d'un vote valide, tout lien $s\to t$ vérifie « $s$ ancêtre strict de $t$ »~;
donc chaque $b_i$ est un ancêtre strict de $b_{i+1}$, et $h(b_i)<h(b_{i+1})$~: la suite des
hauteurs est strictement croissante. En particulier tout justifié est un descendant de $r$.

\textbf{(b)} Soit $H=\max\{h(c) : c\in J\}$ (le maximum existe, $J$ étant fini non vide).
Supposons deux justifiés $c_1\ne c_2$ tels que $h(c_1)=h(c_2)=H$. C'est exactement la
situation exclue par la propriété (iv) admise. Donc le maximum est atteint en un unique
n\oe ud, et la valeur renvoyée par \code{choix_fourche} ne dépend pas de l'ordre de parcours
de l'ensemble $J$~: la règle de choix de fourche est bien définie. \hfill$\square$

\emph{Commentaire.} Cette règle est dite \emph{correcte par construction}~: le théorème de
vivacité plausible (non démontré ici) assure qu'il est \emph{toujours} possible de finaliser
un nouveau point de contrôle au-dessus du justifié de plus grande hauteur, sans qu'aucun
validateur honnête n'ait à enfreindre un commandement.

%=======================================================================
\section*{Partie IV --- Détection des infractions et sûreté responsable}
%=======================================================================

\Q On regroupe les votes par validateur, puis on indexe par hauteur de cible.
\begin{lstlisting}
def infraction_I(votes, h):
    par_val = {}
    for (v, s, t) in votes:
        if v in par_val:
            par_val[v].append((s, t))
        else:
            par_val[v] = [(s, t)]
    for v in par_val:
        vus = {}                      # hauteur de cible -> cible deja vue
        for (s, t) in par_val[v]:
            ht = h[t]
            if ht in vus and vus[ht] != t:
                return (v, vus[ht], t)
            vus[ht] = t
    return None
\end{lstlisting}
\textbf{Correction.} Le dictionnaire \code{vus} mémorise, pour chaque hauteur déjà
rencontrée, la cible correspondante. Une infraction à (I) est exactement l'existence de deux
votes \emph{distincts} d'un même validateur avec $h(t_1)=h(t_2)$~; si les cibles sont égales
($t_1=t_2$) mais les sources différentes, il s'agit d'une infraction à (I) également au sens
strict de l'énoncé du protocole. On a choisi ici de ne signaler que le cas $t_1\ne t_2$, qui
est celui qui met en péril la sûreté (deux branches concurrentes)~; le cas $t_1=t_2$,
$s_1\ne s_2$ se détecterait de la même façon en indexant par $t$.

\textbf{Complexité.} Chaque vote est traité une fois, avec un nombre constant d'opérations de
dictionnaire~: $O(m)$ en moyenne, $O(m)$ en espace.

\Q
\begin{lstlisting}
def infraction_II_naif(votes, h):
    par_val = {}
    for (v, s, t) in votes:
        if v in par_val:
            par_val[v].append((s, t))
        else:
            par_val[v] = [(s, t)]
    for v in par_val:
        L = par_val[v]
        for (s1, t1) in L:
            for (s2, t2) in L:
                if h[s1] < h[s2] and h[s2] < h[t2] and h[t2] < h[t1]:
                    return (v, (s1, t1), (s2, t2))
    return None
\end{lstlisting}
\textbf{Complexité.} Pour un validateur $\nu$ possédant $m_\nu$ votes, la double boucle coûte
$\Theta(m_\nu^2)$. Le coût total est $\sum_\nu m_\nu^2 \le \big(\sum_\nu m_\nu\big)^2 = m^2$,
avec égalité lorsque tous les votes proviennent d'un même validateur~: la complexité est
$O(m^2)$, atteinte dans le pire des cas.

\Q Fixons un validateur $\nu$ et notons $L$ la liste de ses votes.

$(\Rightarrow)$ Supposons que $\nu$ enfreigne (II)~: il existe $(s_1,t_1)$ et $(s_2,t_2)$
dans $L$ avec $h(s_1)<h(s_2)<h(t_2)<h(t_1)$. Alors $(s_1,t_1)$ appartient à l'ensemble
$\{(s,t)\in L : h(s)<h(s_2)\}$, qui est donc non vide, et le maximum des $h(t)$ sur cet
ensemble est $\ge h(t_1) > h(t_2)$. Le vote $(s_2,t_2)$ convient.

$(\Leftarrow)$ Réciproquement, supposons qu'il existe $(s_2,t_2)\in L$ tel que
\[ M \;=\; \max\{h(t_1) : (s_1,t_1)\in L,\ h(s_1)<h(s_2)\} \;>\; h(t_2). \]
Soit $(s_1,t_1)$ un vote réalisant ce maximum. On a alors
$h(s_1)<h(s_2)$ (appartenance à l'ensemble), $h(s_2)<h(t_2)$ (car le vote $(s_2,t_2)$ est
\emph{valide}~: $s_2$ est un ancêtre strict de $t_2$), et $h(t_2)<M=h(t_1)$. Les quatre
inégalités $h(s_1)<h(s_2)<h(t_2)<h(t_1)$ sont donc satisfaites, et les deux votes sont
distincts (leurs sources n'ont pas la même hauteur). Le commandement (II) est enfreint.
\hfill$\square$

C'est ici qu'intervient de façon décisive l'hypothèse de validité des votes~: sans elle, la
condition intermédiaire $h(s_2)<h(t_2)$ devrait être testée explicitement.

\Q L'équivalence précédente suggère de trier par $h(s)$ croissant et de maintenir le maximum
courant des $h(t)$. Il faut traiter les votes de \textbf{même} hauteur de source par
\emph{groupes}, car l'inégalité $h(s_1)<h(s_2)$ est stricte.
\begin{lstlisting}
def infraction_II(votes, h):
    par_val = {}
    for (v, s, t) in votes:
        if v in par_val:
            par_val[v].append((s, t))
        else:
            par_val[v] = [(s, t)]
    for v in par_val:
        L = sorted(par_val[v], key=lambda st: h[st[0]])
        n = len(L)
        M = -1           # max des h(t) sur les votes de source STRICTEMENT plus basse
        temoin = None    # un vote realisant ce maximum
        i = 0
        while i < n:
            j = i                              # [i, j[ = groupe de meme h(s)
            while j < n and h[L[j][0]] == h[L[i][0]]:
                j = j + 1
            for k in range(i, j):              # (s2, t2) candidats
                if M > h[L[k][1]]:
                    return (v, temoin, L[k])
            for k in range(i, j):              # on integre le groupe courant
                if h[L[k][1]] > M:
                    M = h[L[k][1]]
                    temoin = L[k]
            i = j
    return None
\end{lstlisting}
\textbf{Invariant de la boucle principale.}
\begin{quote}\itshape
Au début de chaque itération (d'indice $i$), $M$ est le maximum des $h(t_1)$ sur les votes
$(s_1,t_1)$ de $L$ tels que $h(s_1)<h(L[i]_{\text{source}})$, avec la convention
$M=-1$ si cet ensemble est vide, et \code{temoin} est un vote réalisant ce maximum.
\end{quote}
\emph{Initialisation.} Pour $i=0$, aucun vote n'a de source de hauteur strictement inférieure
à la plus petite~: l'ensemble est vide et $M=-1$.

\emph{Hérédité.} La liste étant triée par $h(s)$ croissant, les votes d'indice $<i$ sont
exactement ceux dont la hauteur de source est $\le h(L[i]_{\text{source}})$~; la tranche
$[i,j[$ regroupe tous ceux de hauteur de source \emph{égale}. Après l'exécution de la seconde
boucle interne, $M$ intègre donc précisément les votes de hauteur de source
$\le h(L[i]_{\text{source}})$, c'est-à-dire strictement inférieure à
$h(L[j]_{\text{source}})$~: l'invariant est rétabli pour $i\leftarrow j$.

\emph{Correction.} Grâce à l'invariant, le test \code{M > h[L[k][1]]} de la première boucle
interne est exactement la condition de la question 35 pour le vote candidat $L[k]$. La
fonction renvoie donc un témoin si et seulement si $\nu$ enfreint (II).

\emph{Terminaison.} La quantité $n-i$ est un entier naturel qui décroît strictement à chaque
tour, car $j>i$ (le groupe est non vide).

\textbf{Complexité.} Le regroupement par validateur coûte $O(m)$ en moyenne. Pour chaque
validateur, le tri coûte $O(m_\nu\log m_\nu)$ et le double parcours $O(m_\nu)$. Comme
$\sum_\nu m_\nu\log m_\nu \le m\log m$, la complexité totale est
\[ \boxed{O(m\log m)} \]
contre $O(m^2)$ pour la version naïve. L'espace supplémentaire est $O(m)$.

\emph{Vérification expérimentale.} Cette fonction et \code{infraction_II_naif} ont été
comparées sur $20\,000$ jeux de votes valides tirés au hasard~: elles concordent toujours, et
le témoin renvoyé vérifie bien $h(s_1)<h(s_2)<h(t_2)<h(t_1)$.

\Q Sur $A_0$, $h=[0,1,2,3,2,3,1,2]$.
\begin{itemize}[leftmargin=1.4em,itemsep=2pt]
\item Votes $(0,3)$ et $(1,2)$~: $h(0)=0$, $h(1)=1$, $h(2)=2$, $h(3)=3$, d'où
$0<1<2<3$, c'est-à-dire $h(s_1)<h(s_2)<h(t_2)<h(t_1)$ avec $(s_1,t_1)=(0,3)$ et
$(s_2,t_2)=(1,2)$~: le validateur enfreint le \textbf{commandement (II)} (vote emboîté).
\item Votes $(0,2)$ et $(0,4)$~: $h(2)=h(4)=2$ et $2\ne 4$~: deux votes distincts pour une
même hauteur de cible, le validateur enfreint le \textbf{commandement (I)}. Notons que $2$ et
$4$ sont précisément en conflit (question 13)~: ce validateur soutient simultanément deux
branches concurrentes.
\end{itemize}

\Q On suppose que l'ensemble $\mathcal{F}$ des validateurs fautifs vérifie
$w(\mathcal{F})<T/3$.

\textbf{(i)} Soient $s_1\to t_1$ et $s_2\to t_2$ deux liens de supermajorité \emph{distincts}
(c'est-à-dire $(s_1,t_1)\ne(s_2,t_2)$) et supposons $h(t_1)=h(t_2)$. Notons $E_1$ et $E_2$
les ensembles de validateurs ayant respectivement voté $s_1\to t_1$ et $s_2\to t_2$~: ce sont
deux supermajorités. D'après le lemme d'intersection (question 7),
$w(E_1\cap E_2)\ge T/3$. Or tout $\nu\in E_1\cap E_2$ a publié les deux votes distincts
$\langle\nu,s_1,t_1\rangle$ et $\langle\nu,s_2,t_2\rangle$ avec $h(t_1)=h(t_2)$~: il enfreint
le commandement (I), donc $\nu\in\mathcal{F}$. Ainsi
$T/3\le w(E_1\cap E_2)\le w(\mathcal{F})<T/3$~: contradiction. Donc $h(t_1)\ne h(t_2)$.
\hfill$\square$

\textbf{(ii)} Même raisonnement. Si $h(s_1)<h(s_2)<h(t_2)<h(t_1)$, les deux liens sont
nécessairement distincts, et tout $\nu\in E_1\cap E_2$ a publié deux votes réalisant
exactement la configuration interdite par le commandement (II). On obtient de nouveau
$w(\mathcal{F})\ge T/3$, contradiction. \hfill$\square$

\Q \textbf{(iii)} Supposons deux liens de supermajorité $s_1\to t_1$ et $s_2\to t_2$ avec
$h(t_1)=h(t_2)=n$. S'ils sont distincts, (i) est contredit. Ils sont donc égaux~: il existe
au plus un lien de supermajorité de hauteur de cible $n$. \hfill$\square$

\textbf{(iv)} Si $n=0$, le seul point de contrôle de hauteur $0$ est la racine. Si $n\ge 1$,
soient $c_1$ et $c_2$ deux justifiés de hauteur $n$. Comme $n\ge 1$, aucun n'est la racine,
donc il existe des liens $c_1'\to c_1$ et $c_2'\to c_2$. Leurs cibles ont même hauteur $n$~;
par (iii) ces deux liens sont égaux, donc en particulier $c_1=c_2$. Il existe donc au plus un
point de contrôle justifié de hauteur $n$. \hfill$\square$

\Q \textbf{(Théorème de sûreté responsable.)} Soient $a_m$ et $b_n$ deux points de contrôle
finalisés en conflit, de fils directs justifiés $a_{m+1}$ et $b_{n+1}$, avec
$h(a_{m+1})=h(a_m)+1$ et $h(b_{n+1})=h(b_n)+1$. Supposons, comme le permet la symétrie,
$h(a_m)<h(b_n)$ (le cas d'égalité est traité à la question suivante), et raisonnons par
l'absurde en supposant $w(\mathcal{F})<T/3$, de sorte que (i)--(iv) s'appliquent.

D'après la question 32(a), il existe une chaîne de justification
\[ r = b_{i_0}\to b_{i_1}\to\dots\to b_{i_k}=b_n \]
dont tous les termes sont justifiés et dont les hauteurs sont strictement croissantes.
Notons pour alléger $\beta_0=r,\ \beta_1,\dots,\beta_k=b_n$ ces points de contrôle.

\emph{Étape 1.} Aucune hauteur $h(\beta_p)$ ne vaut $h(a_m)$ ni $h(a_{m+1})$. En effet,
$\beta_p$ est justifié~; $a_m$ et $a_{m+1}$ le sont aussi ($a_m$ car finalisé, $a_{m+1}$ par
la question 29). Si $h(\beta_p)=h(a_m)$, la propriété (iv) impose $\beta_p=a_m$, donc $a_m$
serait un ancêtre de $b_n$ (les $\beta$ formant une chaîne d'ancêtres de $b_n$), ce qui
contredit le conflit entre $a_m$ et $b_n$. Même argument pour $a_{m+1}$, dont $a_m$ est un
ancêtre.

\emph{Étape 2.} Comme $h(\beta_0)=0\le h(a_m)$ (et $h(\beta_0)\ne h(a_{m+1})$) tandis que
$h(\beta_k)=h(b_n)>h(a_m)$, l'ensemble $\{p : h(\beta_p)>h(a_{m+1})\}$ est non vide dès que
$h(b_n)>h(a_{m+1})$~; c'est le cas, car $h(b_n)>h(a_m)$ et $h(b_n)\ne h(a_{m+1})$ par
l'étape 1. Soit $j$ le plus petit tel indice. Alors $j\ge 1$ et, par minimalité,
$h(\beta_{j-1})<h(a_{m+1})$~; l'étape 1 interdit l'égalité, donc
$h(\beta_{j-1})<h(a_{m+1})$, et même $h(\beta_{j-1})<h(a_m)$ puisque $h(\beta_{j-1})$ ne peut
valoir ni $h(a_m)$ ni $h(a_{m+1})$ et que $h(a_{m+1})=h(a_m)+1$.

\emph{Étape 3.} On dispose donc du lien de supermajorité $\beta_{j-1}\to\beta_j$ avec
\[ h(\beta_{j-1})\;<\;h(a_m)\;<\;h(a_{m+1})\;<\;h(\beta_j), \]
et du lien de supermajorité $a_m\to a_{m+1}$ (qui existe puisque $a_m$ est finalisé). C'est
exactement la configuration interdite par la propriété (ii), avec $(s_1,t_1)=(\beta_{j-1},
\beta_j)$ et $(s_2,t_2)=(a_m,a_{m+1})$~: contradiction.

\emph{Conclusion.} L'hypothèse $w(\mathcal{F})<T/3$ est absurde. Autrement dit, si deux
points de contrôle en conflit sont tous deux finalisés, alors
\[ w(\mathcal{F})\ \ge\ \frac{T}{3}\,, \]
c'est-à-dire qu'au moins un tiers du dépôt total est détruit, et l'on sait \emph{quels}
validateurs punir~: c'est la \textbf{sûreté responsable} (\emph{accountable safety}).
\hfill$\square$

\emph{Commentaire.} Cette propriété est plus forte que la sûreté des algorithmes BFT
classiques~: non seulement une attaque est impossible tant que moins d'un tiers des
participants trichent, mais si elle a lieu, elle est \emph{prouvable} et \emph{sanctionnable}
sur la chaîne. C'est ce qui résout le problème du \emph{nothing at stake} propre à la preuve
d'enjeu.

\Q Si $h(a_m)=h(b_n)$ avec $a_m\ne b_n$ (ils sont en conflit donc distincts), on a deux
points de contrôle justifiés distincts de même hauteur, ce qui contredit directement (iv). En
remontant à la source, ce sont les deux liens $a_{m-1}'\to a_m$ et $b_{n-1}'\to b_n$ dont les
cibles ont même hauteur~: les validateurs de l'intersection (de poids $\ge T/3$) ont publié
deux votes distincts de même hauteur de cible et enfreignent donc le \textbf{commandement
(I)}.

\Q Les tirages sont indépendants, uniformes et avec remise~: à chaque tirage, la probabilité
que le validateur obtenu soit fautif vaut $q$ (proportion de fautifs), indépendamment des
tirages précédents. La variable $X$ compte le rang du premier succès dans une suite d'épreuves
de Bernoulli indépendantes de paramètre $q$~: c'est par définition une \textbf{loi géométrique}
de paramètre $q$, à valeurs dans $\mathbb{N}^*$. Pour $k\ge 1$~:
\[ P(X=k) = (1-q)^{k-1}q,\qquad
   P(X>k) = \sum_{i>k}(1-q)^{i-1}q = (1-q)^{k},\qquad
   \mathbb{E}[X] = \frac{1}{q}. \]
(Le calcul de $P(X>k)$ s'obtient plus simplement en remarquant que $X>k$ signifie que les $k$
premiers tirages sont des échecs, événements indépendants de probabilité $1-q$ chacun.)

\Q L'événement $\{X=+\infty\}$ (l'algorithme ne s'arrête jamais) est l'intersection
décroissante des événements $\{X>k\}$, donc par continuité décroissante de la probabilité
\[ P(X=+\infty) = \lim_{k\to+\infty} P(X>k) = \lim_{k\to+\infty}(1-q)^{k} = 0, \]
puisque $0<1-q<1$. L'algorithme termine donc \textbf{presque sûrement}. \hfill$\square$

\Q
\begin{lstlisting}
import random

def recherche_fautif(validateurs, est_fautif):
    N = len(validateurs)
    n_essais = 0
    while True:
        i = random.randrange(N)
        n_essais = n_essais + 1
        if est_fautif(validateurs[i]):
            return (validateurs[i], n_essais)
\end{lstlisting}
\textbf{Espérance.} Le nombre d'appels à \code{est_fautif} est exactement $X$. Pour
$q=\tfrac13$, $\mathbb{E}[X]=1/q=\mathbf{3}$~: en moyenne trois vérifications suffisent, quel
que soit le nombre $N$ de validateurs. C'est ce qui rend la vérification par échantillonnage
praticable pour un client léger.

\textbf{Pire des cas.} Il n'existe \emph{aucune} borne~: pour tout $k$, $P(X>k)=(1-q)^k>0$.
L'algorithme est de type \emph{Las Vegas}~: son résultat est toujours correct, mais son temps
d'exécution n'est borné qu'en espérance. On ne peut donc parler que de complexité
\emph{en moyenne}, ici $\Theta(1/q)$ appels.

%=======================================================================
\section*{Partie V --- Coût minimal d'une attaque (programmation dynamique)}
%=======================================================================

\Q Les dépôts étant entiers, $w(E)$ est un entier. Pour un entier $w$~:
\[ w\ \ge\ \frac{2T}{3} \iff w\ \ge\ \left\lceil \frac{2T}{3}\right\rceil, \]
car le plus petit entier supérieur ou égal à un réel $x$ est $\lceil x\rceil$. Donc
$C=\lceil 2T/3\rceil$ convient, et c'est bien le plus petit tel entier (pour $C'<C$, l'entier
$w=C'$ vérifierait $w\ge C'$ sans être une supermajorité, dès lors que $C'$ est atteignable).
En Python, avec des divisions entières uniquement~:
\begin{lstlisting}
C = (2 * T + 2) // 3        # ou, de facon equivalente : C = -((-2 * T) // 3)
\end{lstlisting}
\emph{Vérification~:} $T=60\Rightarrow C=40$~; $T=100\Rightarrow C=67$~; $T=101\Rightarrow
C=68$.

\Q \textbf{Appartenance à NP.} Un \emph{certificat} pour l'instance $(d,C)$ est la liste des
indices d'un sous-ensemble $S\subseteq\{0,\dots,n-1\}$. Sa taille est $O(n\log n)$ bits, donc
polynomiale en la taille de l'entrée. Le \emph{vérificateur} calcule $\sum_{i\in S}d_i$ par
$|S|-1\le n-1$ additions d'entiers de taille au plus celle de l'entrée, puis compare à $C$~:
c'est un calcul en temps polynomial. Le problème est donc dans \textbf{NP}. \hfill$\square$

\textbf{Conséquence pour \textsc{Coût-Attaque}.} Un algorithme calculant $\mu(d,C)$ en temps
polynomial résoudrait \textsc{Somme-de-Sous-Ensemble} en temps polynomial~: il suffit de
tester si $\mu(d,C)=C$. En effet, $\mu(d,C)$ est la plus petite somme $\ge C$~; elle vaut
exactement $C$ si et seulement s'il existe un sous-ensemble de somme exactement $C$. Comme
\textsc{Somme-de-Sous-Ensemble} est NP-complet, un tel algorithme entraînerait
$\mathrm{P}=\mathrm{NP}$. Sous l'hypothèse $\mathrm{P}\ne\mathrm{NP}$, il n'existe donc
\textbf{pas} d'algorithme polynomial pour \textsc{Coût-Attaque}.

\Q \textbf{Initialisation.} Pour $i=0$ le seul sous-ensemble de $\varnothing$ est
$\varnothing$, de somme $0$~:
\[ m_0(0)=0 \quad\text{et}\quad m_0(j)=+\infty \ \text{ pour } 1\le j\le C. \]

\textbf{Récurrence.} Un sous-ensemble $S\subseteq\{0,\dots,i\}$ contient ou non l'indice $i$.
En notant $\pi(x)=\min(C,x)$ l'écrêtage, on obtient pour tout $0\le j\le C$~:
\[ m_{i+1}(j) \;=\; \min\Big(\ m_i(j)\ ,\
\min\big\{\,m_i(j')+d_i \ :\ 0\le j'\le C,\ \pi(j'+d_i)=j\,\big\}\ \Big), \]
avec la convention $\min\varnothing=+\infty$. En pratique on implémente cette relation
« en avant »~: pour chaque $j'$ tel que $m_i(j')<+\infty$, on met à jour l'entrée d'indice
$\pi(j'+d_i)$ avec la valeur candidate $m_i(j')+d_i$.

\textbf{Rôle de l'écrêtage.} Sans écrêtage, l'indice $j$ devrait parcourir $[\![0,\sum d_i]\!]$,
tableau dont la taille dépend des données et non du seuil. L'écrêtage borne la taille à $C+1$.
Il ne fait perdre aucune solution optimale~: toutes les sommes $\ge C$ sont regroupées dans la
case $C$, mais on y conserve le \emph{minimum} des valeurs réelles correspondantes, et c'est
précisément ce minimum que l'on cherche. Aucune information utile n'est donc perdue.

\emph{Remarque.} Pour $j<C$ la définition impose $\pi(\sum_{k\in S}d_k)=j$ avec $j<C$, donc
$\sum_{k\in S}d_k = j$~: ainsi $m_i(j)\in\{j,+\infty\}$. Le tableau se comporte comme un
simple test d'atteignabilité pour $j<C$, et seule la case $C$ porte réellement une
optimisation.

\Q Par définition, $\mu(d,C)$ est le minimum des sommes $\ge C$ sur les sous-ensembles de
$\{0,\dots,n-1\}$~; ce sont exactement les sous-ensembles dont la somme écrêtée vaut $C$
(pour $C\ge 1$). Donc
\[ \boxed{\ \mu(d,C) \;=\; m_n(C)\ } \]
(avec la convention $\mu=+\infty$ si $m_n(C)=+\infty$~; et si $C=0$, $\mu=0$).

\Q
\begin{lstlisting}
def cout_min_attaque(d, C):
    INF = float('inf')
    m = [INF] * (C + 1)
    m[0] = 0
    for i in range(len(d)):
        for j in range(C, -1, -1):        # parcours DECROISSANT : indispensable
            if m[j] < INF:
                jp = min(C, j + d[i])
                if m[j] + d[i] < m[jp]:
                    m[jp] = m[j] + d[i]
    return m[C]
\end{lstlisting}
\textbf{Sens de parcours.} Comme $d_i\ge 1$, l'indice cible $j'=\min(C,j+d_i)$ vérifie
$j'>j$, sauf lorsque $j=C$ (où $j'=C$, mais l'affectation est alors sans effet puisque
$m[C]+d_i>m[C]$). En parcourant $j$ de $C$ vers $0$, toute écriture se fait à un indice
$j'>j$ \emph{déjà traité} lors de cette itération~: la valeur écrite ne pourra plus servir de
source dans la même itération. Chaque dépôt est donc utilisé \textbf{au plus une fois}, ce
qui est exactement la sémantique voulue (chaque validateur ne peut être recruté qu'une fois).

\textbf{Contre-exemple en parcours croissant.} Prenons $d=[3]$ et $C=8$. En parcourant $j$ de
$0$ vers $C$~: $j=0$ donne $m[3]=3$~; puis, dans la \emph{même} itération, $j=3$ (dont la
valeur vient d'être écrite) donne $m[6]=6$~; puis $j=6$ donne $m[8]=9$. L'algorithme
renverrait $9$ en utilisant \emph{trois fois} le dépôt $3$, alors que la réponse correcte est
$+\infty$ (un seul dépôt de $3$ ne peut pas atteindre $8$).

\Q \textbf{Invariant.}
\begin{quote}\itshape
À la fin de la $i$-ème itération de la boucle externe ($1\le i\le n$), pour tout
$0\le j\le C$~: $\texttt{m[j]} = m_i(j)$, c'est-à-dire la plus petite somme réalisable à
l'aide d'un sous-ensemble de $\{0,\dots,i-1\}$ dont la somme écrêtée à $C$ vaut exactement
$j$ ($+\infty$ si aucun).
\end{quote}
\emph{Initialisation.} Avant la première itération, $\texttt{m}=[0,+\infty,\dots,+\infty]$,
c'est bien $m_0$.

\emph{Hérédité.} Supposons $\texttt{m}=m_i$ au début de l'itération $i+1$ (indice \code{i}
dans le code). Notons $\texttt{m}^{\text{fin}}$ le tableau à la fin de cette itération.

$(\le)$ Toute valeur écrite dans $\texttt{m}^{\text{fin}}[j']$ est de la forme
$m_i(j)+d_i$ avec $\pi(j+d_i)=j'$, donc réalisable par un sous-ensemble de
$\{0,\dots,i-1\}$ auquel on ajoute $i$ (l'indice $i$ n'était pas dans ce sous-ensemble par
hypothèse de récurrence). Les valeurs conservées sont les $m_i(j)$, réalisables sans $i$.
Donc $\texttt{m}^{\text{fin}}[j']\ge m_{i+1}(j')$ ne peut être violé~: toute valeur du tableau
est réalisable, d'où $\texttt{m}^{\text{fin}}[j']\ge m_{i+1}(j')$.

$(\ge)$ Réciproquement, soit $S\subseteq\{0,\dots,i\}$ réalisant $m_{i+1}(j')$. Si $i\notin
S$, alors $S\subseteq\{0,\dots,i-1\}$ et $m_i(j')\le\sum_{k\in S}d_k$~; comme la case $j'$
n'est jamais \emph{augmentée}, $\texttt{m}^{\text{fin}}[j']\le m_i(j')\le\sum_{k\in S}d_k$.
Si $i\in S$, posons $S'=S\setminus\{i\}$ et $j=\pi\big(\sum_{k\in S'}d_k\big)$. Alors
$m_i(j)\le\sum_{k\in S'}d_k$, et l'itération considère le couple $(j,j'')$ avec
$j''=\pi(j+d_i)$. Deux cas~: si $\sum_{k\in S'}d_k<C$ alors $j=\sum_{k\in S'}d_k$ et
$j''=\pi\big(\sum_{k\in S}d_k\big)=j'$~; si $\sum_{k\in S'}d_k\ge C$ alors $j=C$ et
$j''=C=j'$. Dans les deux cas la mise à jour porte sur la case $j'$ avec la valeur
$m_i(j)+d_i\le \sum_{k\in S}d_k$, d'où $\texttt{m}^{\text{fin}}[j']\le\sum_{k\in S}d_k =
m_{i+1}(j')$.

Enfin, le parcours décroissant garantit (voir question 49) qu'aucune case écrite au cours de
cette itération n'est réutilisée comme source, donc l'indice $i$ n'est jamais compté deux
fois~: les valeurs produites sont bien réalisables sans répétition. L'invariant est établi.
\hfill$\square$

\emph{Conclusion.} À la fin, $\texttt{m[C]}=m_n(C)=\mu(d,C)$ d'après la question 48.

\emph{Terminaison.} Les deux boucles sont des boucles \code{for} sur des intervalles finis~:
la terminaison est immédiate. Le nombre total d'itérations est exactement $n(C+1)$.

\Q \textbf{Complexité.} Temps~: $\Theta\big(n\,C\big)$ (deux boucles imbriquées, corps en
$O(1)$). Espace~: $\Theta(C)$ auxiliaire (un unique tableau de taille $C+1$), contre
$\Theta(nC)$ pour une version conservant tous les $m_i$.

\textbf{Pseudo-polynomialité.} La taille de l'entrée n'est pas $C$ mais le nombre de
\emph{bits} nécessaires pour l'écrire, soit $\Theta(\log C)$. La complexité $\Theta(nC) =
\Theta\big(n\,2^{\log_2 C}\big)$ est donc \emph{exponentielle} en la taille de l'entrée~: on
parle d'algorithme \textbf{pseudo-polynomial}, polynomial en la \emph{valeur} des données mais
non en leur \emph{taille}. Il n'y a donc aucune contradiction avec la NP-complétude établie à
la question 46.

\textbf{Passage aux wei.} Si les dépôts sont exprimés en wei, $C$ est multiplié par $10^{18}$.
La complexité devient de l'ordre de $n\times 10^{18}\times(\text{seuil en ether})$~: le
calcul est totalement irréalisable, et le tableau ne tiendrait pas en mémoire. En pratique on
change d'unité (par exemple le \emph{gwei}, ou mieux le PGCD des dépôts) ou l'on renonce à
l'optimum exact au profit d'une heuristique gloutonne (trier les dépôts par ordre décroissant
et recruter jusqu'à atteindre $C$), qui donne ici une $2$-approximation.

\Q Avec $d=[12,\,7,\,20,\,5,\,16]$~:
\[ T = 12+7+20+5+16 = 60,\qquad C = \left\lceil \frac{2\times 60}{3}\right\rceil = 40. \]
On trouve $\mu(d,40)=\mathbf{40}$, atteint par le sous-ensemble
\[ S^\star=\{12,\ 7,\ 5,\ 16\},\qquad 12+7+5+16=40. \]
(C'est le seul sous-ensemble de somme exactement $40$.) L'attaquant doit donc contrôler
quatre validateurs sur cinq, mais peut se dispenser du plus gros dépôt.

\textbf{Déroulé du tableau} pour la version réduite $d=[3,5,7]$ et $C=8$ (on note
$\infty$ pour $+\infty$)~:
\[
\begin{array}{l|ccccccccc}
j & 0&1&2&3&4&5&6&7&8\\\hline
\text{initialisation} & 0&\infty&\infty&\infty&\infty&\infty&\infty&\infty&\infty\\
\text{après } d_0=3 & 0&\infty&\infty&3&\infty&\infty&\infty&\infty&\infty\\
\text{après } d_1=5 & 0&\infty&\infty&3&\infty&5&\infty&\infty&8\\
\text{après } d_2=7 & 0&\infty&\infty&3&\infty&5&\infty&7&8
\end{array}
\]
D'où $\mu([3,5,7],8) = m_3(8) = \mathbf{8}$, réalisé par $\{3,5\}$. On vérifie au passage la
remarque de la question 47~: pour $j<8$, chaque case vaut $j$ ou $\infty$~; seule la case
$j=8$ agrège plusieurs sommes ($8$, $10$, $12$, $15$) et n'en garde que le minimum $8$.
Noter aussi que la case $j=8$ passe de $\infty$ à $8$ dès le traitement de $d_1=5$ (via
$3+5$) et n'est plus améliorée ensuite.

\Q Sur l'exemple, $\mu(d,C)=40=C$~: l'attaquant atteint \emph{exactement} le seuil.

\textbf{Ce n'est pas toujours possible.} Prenons $d=[1,\,100]$. Alors $T=101$ et
$C=\lceil 202/3\rceil = 68$. Les sommes réalisables sont $0$, $1$, $100$ et $101$~; les
seules $\ge 68$ sont $100$ et $101$, donc
\[ \mu(d,C) = 100 \;>\; 68 = C. \]
La granularité des dépôts empêche d'atteindre le seuil~: pour réunir une supermajorité,
l'attaquant est contraint de mobiliser $100$, soit près de $99\,\%$ du dépôt total, alors que
$67{,}3\,\%$ auraient « suffi » en théorie. Un système où quelques validateurs concentrent
l'essentiel du dépôt rend donc l'attaque \emph{plus chère en valeur absolue}, mais aussi
beaucoup plus facile à monter en pratique (il suffit de corrompre un seul acteur)~: la
métrique pertinente pour la sécurité réelle est le nombre d'acteurs à corrompre autant que le
capital à mobiliser.

\textbf{Lien avec la question 7.} Il faut distinguer deux quantités~:
\begin{itemize}[leftmargin=1.4em,itemsep=1pt]
\item le dépôt \emph{mobilisé} pour former une supermajorité, au moins
$C=\lceil 2T/3\rceil$ (ici $40$, soit $\tfrac23 T$)~;
\item le dépôt effectivement \emph{détruit} lors d'une violation de la sûreté, au moins
$T/3$ (ici $20$) d'après le lemme d'intersection.
\end{itemize}
Un attaquant qui contrôle $2T/3$ et finalise deux points de contrôle en conflit perd donc au
minimum la moitié de ce qu'il a mobilisé. C'est ce ratio, et non la seule difficulté
algorithmique, qui constitue la véritable barrière économique de Casper.

\Q \textbf{Synthèse.} Les parties III et IV établissent la \emph{sûreté responsable}~: sous
l'hypothèse que moins d'un tiers du dépôt triche, deux points de contrôle en conflit ne
peuvent être finalisés, et si l'hypothèse tombe, les coupables sont identifiables et punis
(partie IV, question 40). Cette garantie est \emph{inconditionnelle vis-à-vis du réseau}~:
elle ne suppose aucune borne sur la latence.

En revanche elle ne dit rien de la \emph{vivacité}. Si plus d'un tiers des validateurs
cessent de voter — panne massive, partition réseau, censure — plus aucune supermajorité ne
peut se former~: aucun lien de supermajorité n'apparaît, l'ensemble $J$ de la question 26
cesse de croître, et la chaîne ne finalise plus rien. Le protocole est bloqué sans que
personne n'ait triché.

C'est exactement le rôle de la \textbf{fuite d'inactivité} de la partie I~: en faisant
décroître géométriquement le dépôt des absents (question 8), elle diminue le dépôt total $T$
et donc le seuil $C=\lceil 2T/3\rceil$, jusqu'à ce que les validateurs restés actifs
redeviennent, à eux seuls, une supermajorité. Le nombre d'époques nécessaires est le
$k^\star$ de la question 9, logarithmique en le rapport des dépôts~: la reprise est
relativement rapide. La vivacité est ainsi restaurée.

Mais ce mécanisme a un prix. La démonstration de la question 40 repose sur le fait que les
deux supermajorités en cause sont mesurées \emph{avec le même} dépôt total $T$~; c'est ce qui
autorise le lemme d'intersection. Si le réseau se scinde en deux partitions $\mathcal{V}_A$ et
$\mathcal{V}_B$, chacune voit l'autre comme inactive et lui applique la fuite~: sur la chaîne
$A$ le dépôt de $\mathcal{V}_B$ fond, et réciproquement. Au bout d'un certain temps,
$\mathcal{V}_A$ est une supermajorité \emph{selon les comptes de la chaîne $A$} et
$\mathcal{V}_B$ une supermajorité \emph{selon ceux de la chaîne $B$}~: deux points de contrôle
en conflit sont finalisés sans qu'aucun validateur n'ait enfreint le moindre commandement,
donc \textbf{sans qu'aucun dépôt ne soit détruit}. L'hypothèse du lemme d'intersection est
violée, non par malveillance, mais parce que les deux supermajorités ne se rapportent plus au
même $T$.

Casper échange donc une garantie de sûreté \emph{absolue} contre une garantie de sûreté
\emph{sous hypothèse de synchronie faible}, assortie d'une garantie de vivacité. C'est un
compromis classique et inévitable~: le théorème CAP, comme l'impossibilité FLP, interdit
d'obtenir simultanément la sûreté et la vivacité dans un réseau asynchrone sujet aux
partitions. Le protocole prescrit alors une règle simple et purement locale~: \emph{chaque
validateur retient le premier point de contrôle finalisé qu'il a vu}, et l'arbitrage entre les
deux branches est renvoyé hors du protocole, à un embranchement logiciel décidé par la
communauté.

\vfill
\begin{center}\small\textsc{Fin du corrigé} \quad---\quad
\href{https://cpgeacademy.org}{cpgeacademy.org}\end{center}
\end{document}
